\documentclass[10pt,a4paper]{article}
\usepackage[latin1]{inputenc}
\usepackage[T1]{fontenc}
\usepackage[french]{babel}
\usepackage{fancyhdr}
\usepackage{graphicx}
\pagestyle{fancy}

\renewcommand{\headrulewidth}{1pt}
%\fancyhead[C]{\textbf{page \thepage}}
\fancyhead[L]{\includegraphics[width=2cm]{../images/enib}}
\fancyhead[R]{S6A-02}

\renewcommand{\footrulewidth}{1pt}
%\fancyfoot[L]{\textbf{page \thepage}}
\fancyfoot[C]{Corentin Pouplard - Fidel Nguyen - Simon Tarti�re -
  Yoann Diqu�lou}
%\fancyfoot[R]{\leftmark}

\title{{\Huge \bf Mod�lisation du pont de recouvrance}\\
  \vspace{0.5cm}
\centerline{  Corentin Pouplard - Fidel Nguyen}
\centerline{ Simon Tarti�re -  Yoann Diqu�lou}
\vspace{0.5cm}
\date{\today}
}
\begin{document}
\maketitle
\thispagestyle{fancy}
\hrule
\section*{Objectif}
Le but de ce projet est de mod�liser le \textit{Pont national de
  Brest} tel qu'il �tait de 1861 � 1944.\\
Nous souhaitons permettre la visualisation de ses m�canismes et les
observer en fonctionnement tout en s'assurant de la coh�rence de nos mod�les.\\
\section*{Fonctionnalit�s}
Les fonctionnalit�s pr�vues pour notre application sont:
\begin{itemize}
\item Mettre en mouvement le pont.
\item Se d�placer sur le pont.
\item Pouvoir consulter l'historique du pont.
\end{itemize}

\section*{Perspectives d'am�lioration}
Voici quelques am�liorations que nous pourions mettre en place si le
temps nous le permet:
\begin{itemize}
\item Ajout d'intelligences artificielles (IA).
\item Faire passer un bateau sous le pont.
\item Possibilit� de vue �clat�e.
\item Commander les IA.
\end{itemize}
\end{document}
